\documentclass[a4paper, 12pt]{article}
\usepackage{graphicx} % Required for inserting images
\usepackage{fullpage}
\usepackage{amsmath}
\usepackage{xcolor}
\usepackage{float}
\usepackage{geometry}
\usepackage{biblatex}
\geometry{margin=1in}
\usepackage{enumitem}
\usepackage{hyperref}
\usepackage{microtype}
\usepackage{parskip}

\hypersetup{
    colorlinks=true,        % Enable colored links
    linkcolor=teal,         % Set color for internal links
    citecolor=teal,         % Set color for citations
    filecolor=teal,         % Set color for file links
    urlcolor=teal           % Set color for URLs
}

\title{01-11-2026 Decimals}
\author{C\&L Math Tutoring}
\date{January 2025}

\begin{document}

\maketitle

\subsection*{Decimal place values}

Since we've already covered fractions, we can now explore decimals, which are another way to represent parts of a whole.

$$3,567,432.245895$$

Numbers after the decimal point represent parts of a whole, just like
fractions.

Let's look at place value:

Three \textbf{millions} $10^6$\\
Five \textbf{hundred thousands} $10^5$\\
Six \textbf{ten thousands} $10^4$\\
Seven \textbf{thousands} $10^3$\\
Four \textbf{hundreds} $10^2$\\
Three \textbf{tens} $10^1$\\
Two \textbf{ones} $10^0$\\
\textbf{Decimal point goes here} \\
Two \textbf{tenths} $10^{-1}$\\
Four \textbf{hundredths} $10^{-2}$ \\
Five \textbf{thousandths} $10^{-3}$ \\
Eight \textbf{ten thousandths} $10^{-4}$ \\
Nine \textbf{hundred thousandths} $10^{-5}$ \\
Five \textbf{millionths}$10^{-6}$

\subsection*{Word form}

To get the word form of a decimal:

\begin{enumerate}
\item Say the whole number part normally.
\item Say ``and" for the decimal point.
\item For the decimal part, say it as though it were a whole number, then say the place value of the last digit.
\end{enumerate}

Three million, five hundred sixty-seven thousand, four hundred thirty-two and two hundred forty-five thousand, eight hundred ninety-five millionths.

\subsection*{Conversions}

Four \textit{hundredths} = $\displaystyle \frac{4}{100} = \frac{1}{25} = 0.04$

To convert from a fraction to a decimal, perform long division of the numerator by the denominator.

To convert from a decimal to a fraction, write the decimal as a fraction with a denominator that is a power of ten (last place value), then simplify.

\subsection*{Adding and subtracting decimals}

When adding and subtracting decimals, line up the decimal point and add/subtract as usual.

Example 1: 0.75 + 2.4

\begin{align*}
0&.75 \\
+ \: 2&.40 \\
= \: 3&.15
\end{align*}

Example 2: 5.2 - 3.75

\begin{align*}
5&.20 \\
- \: 3&.75 \\
= \: 1&.45
\end{align*}

\subsection*{Multiplying and dividing decimals}

When multiplying decimals, multiply as usual, then place the decimal point in the product so that it has the same number of decimal places as the sum of the decimal places in the factors.

\end{document}